\section{Task and Design Requirements}
\label{sec:setting}

The task is to change an object's planar pose relative to a parallel gripper while retaining it in the grasp. The gripper pinches the object across its thickness, with gravity normal to the manipulation plane, so the object's weight acts along the pinch axis and does not enter the planar analysis. At engagement the gripper's world pose is $(\mathbf r^w_0,\psi_0)$; the wrist orientation $\psi_0$ is then held fixed. Vectors are expressed in the object frame, which coincides with the gripper frame at engagement, has its origin at the initial grasp, and moves with the object. A superscript $w$ marks the same vector in the world frame; for example, a rod centre is $\mathbf q$ in the object frame and $\mathbf q^w$ in the world frame. A request specifies the object's final pose relative to the gripper: a turn $\Delta$ and a displacement $\mathbf d$, with $\mathbf Q(\cdot)$ the planar rotation by a given angle. We take $\Delta\geq0$, a counterclockwise turn, without loss of generality: a clockwise turn is the mirror image of the same construction.

The object is rigid, with a polygonal outline $P$, convex or not, and a known initial grasp in its interior. Manipulation is quasi-static. The gripper touches the object through a pad centred on the grasp, and the grasp is usable while the pad's footprint lies inside $P$; these grasp positions form the usable grasp region. The grip must retain the object while permitting sliding and pivoting when the object is pressed against the tool. Planning uses no friction coefficient or grip-force value. The only friction model we invoke, in Section~\ref{sec:slip}, is the shape of the pad's limit surface, and the coefficient and force cancel from it.

The design objective is a passive workspace tool that enables both changes of position and changes of orientation without adding finger actuation or requiring wrist rotation during manipulation. Its geometry must supply the contact constraints, and its placement must admit robot access, usable grasp travel, and collision-free motion.
